\include{./def/yf-formatting}
\include{./def/yf-def}

\title{}
\author{}
\date{}



\def\bD{\mathbb{D}}
\def\bQ{\mathbb{Q}}

\def\A{\mathcal{A}}
\def\B{\mathcal{B}}
\def\C{\mathcal{C}}
\def\D{\mathcal{D}}
\def\E{\mathcal{E}}
\def\F{\mathcal{F}}
\def\G{\mathcal{G}}
\def\I{\mathcal{I}}
\def\L{\mathcal{L}}
\def\R{\mathcal{R}}
\def\T{\mathcal{T}}
%\def\U{\mathcal{U}}
%\def\X{\mathcal{X}}
\def\rank{\mit{rank}}
\def\tm{\mit{Time}}

\def\extraspacing{\vspace{3mm} \noindent}
\def\vgap{\vspace{3mm}}
\def\vslit{\vspace{0.5mm}}

\DeclareMathOperator*{\Log}{Log}

\begin{document}

\boxminipg{\linewidth}{
    \begin{center}
        \vspace{3mm}
    {\Large
    (CE7456 Lecture Notes 9) Subgraph Enumeration} \\[2mm]
    {\large Yufei Tao}
    \vspace{3mm}
    \end{center}
}

Today, we will discuss the {\em subgraph enumeration} problem. Let $G = (V, E)$ be an undirected graph called the {\em data graph}. A graph $G' = (V', E')$ is a {\em subgraph} of $G$ if $V' \subseteq V$ and $E' \subseteq E$. Let $Q = (V_Q, E_Q)$ be another undirected graph called the {\em pattern graph}. A subgraph $G' = (V', E')$ of $G$ is {\em isomorphic} to $Q$ if there is a bijection function $f: V_Q \rightarrow V'$ such that, for any distinct vertices $u, v \in V'$, the pattern graph $Q$ has an edge $\set{u, v}$ if and only if the edge $\set{f(u), f(v)}$ exists in $G'$. The goal of subgraph enumeration is to report all the subgraphs of $G$ that are isomorphic to $Q$. These subgraphs are called the {\em occurrences} of $Q$.

\vgap

We assume that $Q$ has a constant size, i.e., the number $|V_Q|$ of its vertices can be regarded as a constant (and, hence, so can the number $|E_Q|$ of its edges). This is necessary; otherwise, the problem is known to be NP-hard: finding whether $G$ has a clique is NP-hard.

\section{Fractional Edge Cover Number} \label{sec:frac-edge-cover-num}

A function $w: E_Q \rightarrow [0, 1]$ is said to be a {\em fractional edge cover} of $Q = (V_Q, E_Q)$ if, for every vertex $u \in V_Q$, it holds that
\myeqn{
    \sum_{e \in E_Q: u \in e} w(e) &\ge& 1. \label{eqn:frac-edge-cover}
}
The {\em total weight} of $w$ is defined as
\myeqn{
    \sum_{e \in E_Q} w(e). \nn
}
The {\em fractional edge cover number} of $Q$ --- denoted as $\rho(Q)$ --- is defined as the smallest total weight of all fractional edge covers of $Q$.

\vgap

For example, if $Q$ is a $k$-cycle (i.e., a cycle with $k$ edges), then $\rho(Q) = k/2$; if $Q$ is a $k$-star (i.e., a graph with $k + 1$ vertices $u, v_1, ..., v_k$, and $k$ edges $\set{u, v_1}, \set{u, v_2}, ..., \set{u, v_k}$), then $\rho(Q) = k$.

\vgap

The following is known as the {\em AGM bound}:

\begin{theorem} \label{thm:agm}
    If the data graph $G$ has $m$ edges and the pattern graph $Q$ has a constant size, then $G$ can have $O(m^{\rho(Q)})$ occurrences of $Q$.
\end{theorem}

The AGM bound is tight:

\begin{theorem} \label{thm:agm}
    Fix any integer $m \ge 1$ and any pattern graph $Q$ with a constant size. There is a data graph $G$ with at most $m$ edges that has $\Omega(m^{\rho(Q)})$ occurrences of $Q$.
\end{theorem}

We will prove both theorems in this lecture.

\end{document}
